En el capítulo anterior estudiamos el movimiento rectilíneo uniforme, en el
cual la velocidad permanece constante.

Sin embargo, la mayoría de los movimientos que observamos no son así.

Un automóvil puede arrancar desde el reposo, aumentar su rapidez, frenar o
cambiar de sentido. En todas estas situaciones la velocidad está cambiando.

La magnitud física que nos permite describir ese cambio es la
\textbf{aceleración}.

\begin{notapedagogica}
Una forma directa de comenzar la clase es plantear tres situaciones:

\begin{enumerate}
    \item un automóvil pasa de \(\SI{0}{\meter\per\second}\) a
    \(\SI{20}{\meter\per\second}\);

    \item un automóvil pasa de \(\SI{20}{\meter\per\second}\) a
    \(\SI{5}{\meter\per\second}\);

    \item un automóvil que se mueve hacia la izquierda cambia su velocidad
    de \(-\SI{5}{\meter\per\second}\) a
    \(-\SI{15}{\meter\per\second}\).
\end{enumerate}

En los tres casos hay aceleración porque la velocidad cambia.

Esto permite corregir desde el comienzo la idea equivocada de que
``acelerar'' significa únicamente ``ir cada vez más rápido''.
\end{notapedagogica}


\section{El concepto de aceleración}

La velocidad nos indica qué tan rápidamente cambia la posición.

Ahora queremos describir qué tan rápidamente cambia la propia velocidad.

\begin{definicion}[Aceleración]
La \textbf{aceleración} es la magnitud vectorial que describe la razón de
cambio de la velocidad con respecto al tiempo.
\end{definicion}

En movimiento rectilíneo trabajaremos con la componente de aceleración sobre
el eje elegido, que representaremos mediante \(a_x\), o simplemente \(a\)
cuando no exista posibilidad de confusión.

\begin{idea}
La aceleración no nos dice directamente qué tan rápido se mueve un cuerpo.

Nos dice \textbf{qué tan rápidamente está cambiando su velocidad}.
\end{idea}


\section{Aceleración media}

Supongamos que en el instante \(t_i\) una partícula tiene velocidad \(v_i\)
y que en un instante posterior \(t_f\) tiene velocidad \(v_f\).

El cambio de velocidad es

\[ \Delta v=v_f-v_i. \]

El intervalo de tiempo es

\[ \Delta t=t_f-t_i. \]

\begin{definicion}[Aceleración media]
La \textbf{aceleración media} durante un intervalo de tiempo es el cociente
entre el cambio de velocidad y el intervalo de tiempo en el que ocurre.
\end{definicion}

En movimiento rectilíneo,

\[ \boxed{\bar{a}=\frac{\Delta v}{\Delta t}
=\frac{v_f-v_i}{t_f-t_i}}. \]

La unidad de aceleración en el Sistema Internacional es

\[ \si{\meter\per\second\squared}. \]


\subsection{¿Qué significa \(\si{\meter\per\second\squared}\)?}

La unidad

\[ \si{\meter\per\second\squared} \]

puede interpretarse como

\[ \frac{\si{\meter\per\second}}{\si{\second}}. \]

Por ejemplo, si

\[ a=+\SI{3}{\meter\per\second\squared}, \]

significa que la componente de la velocidad aumenta

\[ \SI{3}{\meter\per\second} \]

por cada segundo transcurrido, siempre que la aceleración permanezca
constante.

Si inicialmente

\[ v_0=\SI{2}{\meter\per\second}, \]

entonces:

\[
\begin{array}{c|c}
t\;(\si{\second}) & v\;(\si{\meter\per\second})\\
\hline
0 & 2\\
1 & 5\\
2 & 8\\
3 & 11
\end{array}
\]

\begin{interpretacion}
Una aceleración de

\[ \SI{3}{\meter\per\second\squared} \]

no significa que el cuerpo recorra tres metros cada segundo.

Significa que su \textbf{velocidad cambia}
\(\SI{3}{\meter\per\second}\) cada segundo.
\end{interpretacion}


\begin{ejemplo}[Cálculo de aceleración media]

La velocidad de un automóvil cambia de

\[ v_i=\SI{8}{\meter\per\second} \]

a

\[ v_f=\SI{20}{\meter\per\second} \]

en un intervalo de

\[ \Delta t=\SI{4}{\second}. \]

La aceleración media es

\[
\bar{a}
=
\frac{v_f-v_i}{\Delta t}.
\]

Sustituyendo,

\[
\bar{a}
=
\frac{
\SI{20}{\meter\per\second}
-
\SI{8}{\meter\per\second}
}{
\SI{4}{\second}
}.
\]

Por tanto,

\[ \boxed{\bar{a}=\SI{3}{\meter\per\second\squared}}. \]

Esto significa que, en promedio, la velocidad aumentó
\(\SI{3}{\meter\per\second}\) por cada segundo.
\end{ejemplo}


\section{Aceleración instantánea}

La aceleración media describe el cambio de velocidad durante un intervalo.

Pero, al igual que hicimos con la velocidad, podemos preguntarnos:

\begin{quote}
¿Cuál es la aceleración exactamente en un instante determinado?
\end{quote}

Para responder, hacemos cada vez más pequeño el intervalo de tiempo utilizado
para calcular la aceleración media.

\begin{definicion}[Aceleración instantánea]
La \textbf{aceleración instantánea} es el límite de la aceleración media
cuando el intervalo de tiempo tiende a cero.
\end{definicion}

Matemáticamente,

\[ \boxed{a=\lim_{\Delta t\to0}\frac{\Delta v}{\Delta t}}. \]

Utilizando cálculo diferencial,

\[ \boxed{a=\frac{\dd v}{\dd t}}. \]

Como

\[ v=\frac{\dd x}{\dd t}, \]

también podemos escribir

\[ \boxed{a=\frac{\dd^2x}{\dd t^2}}. \]

\begin{interpretacion}
La aceleración es la tasa instantánea de cambio de la velocidad.

Por eso la relación entre posición, velocidad y aceleración puede resumirse
como

\[
x(t)
\quad
\xrightarrow{\;\deriv{}{t}\;}
\quad
v(t)
\quad
\xrightarrow{\;\deriv{}{t}\;}
\quad
a(t).
\]
\end{interpretacion}


\section{El signo de la aceleración}

En movimiento rectilíneo, el signo de la aceleración depende del sistema de
referencia elegido.

Si escogemos hacia la derecha como sentido positivo:

\begin{itemize}
    \item \(a>0\) significa que la velocidad cambia en el sentido positivo;
    \item \(a<0\) significa que la velocidad cambia en el sentido negativo.
\end{itemize}

\begin{advertencia}
Una aceleración negativa no significa necesariamente que el cuerpo esté
frenando.

El signo de la aceleración indica su dirección respecto del eje elegido.

Para saber si la rapidez aumenta o disminuye debemos comparar los signos de
la velocidad y la aceleración.
\end{advertencia}


\section{¿Cuándo aumenta o disminuye la rapidez?}

Esta relación es fundamental.

Si velocidad y aceleración tienen el \textbf{mismo signo}, la rapidez
aumenta.

Si tienen \textbf{signos opuestos}, la rapidez disminuye.

\begin{table}[H]
    \centering

    \caption{Relación entre los signos de velocidad y aceleración.}

    \begin{tabular}{ccc}
        \toprule
        Velocidad & Aceleración & Comportamiento de la rapidez \\
        \midrule
        \(v>0\) & \(a>0\) & aumenta \\
        \(v>0\) & \(a<0\) & disminuye \\
        \(v<0\) & \(a<0\) & aumenta \\
        \(v<0\) & \(a>0\) & disminuye \\
        \bottomrule
    \end{tabular}
\end{table}

\begin{idea}
La regla correcta es:

\[
\boxed{
v\text{ y }a\text{ con el mismo signo}
\Longrightarrow
\text{aumenta la rapidez}
}
\]

\[
\boxed{
v\text{ y }a\text{ con signos opuestos}
\Longrightarrow
\text{disminuye la rapidez}
}
\]
\end{idea}


\begin{ejemplo}[Velocidad y aceleración negativas]

Un automóvil se mueve hacia la izquierda con

\[ v=-\SI{10}{\meter\per\second} \]

y tiene aceleración

\[ a=-\SI{2}{\meter\per\second\squared}. \]

La velocidad y la aceleración tienen el mismo signo.

Después de un segundo, si la aceleración permanece constante,

\[ v=-\SI{12}{\meter\per\second}. \]

La rapidez pasó de

\[ \SI{10}{\meter\per\second} \]

a

\[ \SI{12}{\meter\per\second}. \]

Por tanto, el automóvil está aumentando su rapidez aunque su aceleración sea
negativa.
\end{ejemplo}


\begin{ejemplo}[Aceleración positiva y disminución de rapidez]

Una partícula se mueve hacia la izquierda con

\[ v=-\SI{12}{\meter\per\second} \]

y tiene aceleración

\[ a=+\SI{3}{\meter\per\second\squared}. \]

La velocidad es negativa y la aceleración positiva.

Tienen signos opuestos.

Por ello la rapidez disminuye.

Un segundo después,

\[ v=-\SI{9}{\meter\per\second}. \]

La partícula continúa moviéndose hacia la izquierda, pero cada vez con menor
rapidez.
\end{ejemplo}


\section{Velocidad nula no significa aceleración nula}

Un cuerpo puede tener velocidad cero en un instante y, sin embargo, tener
aceleración distinta de cero.

Por ejemplo, un objeto puede disminuir su rapidez hasta detenerse
momentáneamente y después comenzar a moverse en el sentido contrario.

En el instante exacto en que cambia de sentido,

\[ v=0, \]

pero la aceleración puede continuar siendo distinta de cero.

\begin{idea}
La velocidad describe el estado instantáneo del movimiento.

La aceleración describe cómo está cambiando ese estado.

Por eso

\[ v=0 \]

no implica necesariamente

\[ a=0. \]
\end{idea}


\section{Interpretación de la gráfica velocidad-tiempo}

Consideremos una gráfica de velocidad en función del tiempo.

La pendiente entre dos puntos es

\[
m
=
\frac{v_f-v_i}{t_f-t_i}.
\]

Pero

\[
v_f-v_i=\Delta v
\]

y

\[
t_f-t_i=\Delta t.
\]

Por tanto,

\[ \boxed{m=\frac{\Delta v}{\Delta t}=\bar{a}}. \]

\begin{idea}
En una gráfica \(v-t\),

\[
\boxed{
\text{pendiente de una secante}
=
\text{aceleración media}
}.
\]

En el límite,

\[
\boxed{
\text{pendiente de la tangente}
=
\text{aceleración instantánea}
}.
\]
\end{idea}


\section{Movimiento con aceleración constante}

Estudiaremos ahora un caso particular muy importante.

\begin{definicion}[Movimiento rectilíneo uniformemente acelerado]
El \textbf{movimiento rectilíneo uniformemente acelerado}, abreviado MRUA,
es un movimiento rectilíneo en el cual la aceleración permanece constante
durante el intervalo considerado.
\end{definicion}

En este movimiento,

\[ a=\text{constante}. \]

Esto significa que la velocidad experimenta cambios iguales durante
intervalos de tiempo iguales.

Por ejemplo, si

\[ a=\SI{2}{\meter\per\second\squared}, \]

la velocidad aumenta

\[ \SI{2}{\meter\per\second} \]

cada segundo.

\begin{nota}
Matemáticamente, el caso \(a=0\) también corresponde a una aceleración
constante.

En ese caso el MRUA se reduce al movimiento rectilíneo uniforme.

Habitualmente, cuando hablamos de MRUA nos interesa especialmente el caso

\[ a\neq0. \]
\end{nota}


\section{Primera ecuación: velocidad en función del tiempo}

Para aceleración constante,

\[ a=\frac{\Delta v}{\Delta t}. \]

Por tanto,

\[ a=\frac{v_f-v_i}{t_f-t_i}. \]

Multiplicando,

\[ v_f-v_i=a(t_f-t_i). \]

Así,

\[ \boxed{v_f=v_i+a(t_f-t_i)}. \]

Si elegimos

\[ t_i=0, \]

y escribimos la velocidad inicial como \(v_0\), obtenemos

\[ \boxed{v=v_0+at}. \]

\begin{interpretacion}
La velocidad final es igual a:

\[
\text{velocidad inicial}
+
\text{cambio de velocidad}.
\]

El término

\[ at \]

representa cuánto cambió la velocidad durante el intervalo.
\end{interpretacion}


\begin{ejemplo}[Velocidad después de cierto tiempo]

Un automóvil tiene velocidad inicial

\[ v_0=\SI{5}{\meter\per\second} \]

y aceleración constante

\[ a=\SI{3}{\meter\per\second\squared}. \]

¿Cuál es su velocidad después de

\[ t=\SI{4}{\second}? \]

Utilizamos

\[ v=v_0+at. \]

Sustituyendo,

\[
v
=
\SI{5}{\meter\per\second}
+
\left(\SI{3}{\meter\per\second\squared}\right)
\left(\SI{4}{\second}\right).
\]

Por tanto,

\[
v
=
\SI{5}{\meter\per\second}
+
\SI{12}{\meter\per\second}.
\]

Así,

\[ \boxed{v=\SI{17}{\meter\per\second}}. \]
\end{ejemplo}


\section{Gráfica velocidad-tiempo en MRUA}

Como

\[ v=v_0+at, \]

la velocidad es una función lineal del tiempo.

Por tanto, la gráfica \(v-t\) es una línea recta.

\begin{figure}[H]
    \centering

    \begin{tikzpicture}
        \begin{axis}[
            axis lines=middle,
            xlabel={$t$},
            ylabel={$v$},
            xmin=0,
            xmax=6,
            ymin=0,
            ymax=16,
            width=9cm,
            height=6cm
        ]

            \addplot[
                domain=0:5,
                samples=2,
                thick
            ] {3+2*x};

        \end{axis}
    \end{tikzpicture}

    \caption{Gráfica velocidad-tiempo para una aceleración constante positiva.}
\end{figure}

La pendiente es

\[ a. \]

El punto donde la recta corta el eje vertical es

\[ v_0. \]

Por tanto,

\[
\boxed{
\text{pendiente de la gráfica }v-t=a
}.
\]


\section{Desplazamiento como área bajo la gráfica velocidad-tiempo}

En el capítulo anterior establecimos que el área algebraica bajo una gráfica
velocidad-tiempo representa el desplazamiento.

Para aceleración constante, la gráfica entre \(v_0\) y \(v\) forma un
trapecio.

Su área es

\[
\Delta x
=
\frac{v_0+v}{2}\,t.
\]

Por tanto,

\[ \boxed{\Delta x=\frac{v_0+v}{2}\,t}. \]

Esto permite introducir una propiedad especial del MRUA.

\begin{definicion}[Velocidad media con aceleración constante]
Cuando la aceleración es constante, la velocidad media durante un intervalo
es igual al promedio aritmético entre la velocidad inicial y la velocidad
final:
\end{definicion}

\[ \boxed{\bar{v}=\frac{v_0+v}{2}}. \]

\begin{advertencia}
La expresión

\[ \bar{v}=\frac{v_0+v}{2} \]

no es válida para cualquier movimiento.

Es válida cuando la velocidad cambia linealmente con el tiempo, es decir,
cuando la aceleración es constante.
\end{advertencia}


\section{Segunda ecuación: posición en función del tiempo}

Partimos de

\[ \Delta x=\frac{v_0+v}{2}\,t. \]

Sabemos también que

\[ v=v_0+at. \]

Sustituimos:

\[
\Delta x
=
\frac{v_0+(v_0+at)}{2}\,t.
\]

Simplificando,

\[
\Delta x
=
\frac{2v_0+at}{2}\,t.
\]

Entonces,

\[ \boxed{\Delta x=v_0t+\frac{1}{2}at^2}. \]

Como

\[ \Delta x=x-x_0, \]

obtenemos

\[ \boxed{x=x_0+v_0t+\frac{1}{2}at^2}. \]

\begin{interpretacion}
La posición tiene tres contribuciones:

\[
x
=
\underbrace{x_0}_{\text{posición inicial}}
+
\underbrace{v_0t}_{\text{movimiento que tendría sin acelerar}}
+
\underbrace{\frac{1}{2}at^2}_{\text{efecto de la aceleración}}.
\]
\end{interpretacion}


\subsection{Comprobación dimensional}

Cada término debe tener unidades de longitud.

Para

\[ v_0t, \]

tenemos

\[
\left(\si{\meter\per\second}\right)
(\si{\second})
=
\si{\meter}.
\]

Para

\[ \frac{1}{2}at^2, \]

tenemos

\[
\left(\si{\meter\per\second\squared}\right)
(\si{\second\squared})
=
\si{\meter}.
\]

Por tanto, la ecuación es dimensionalmente homogénea.


\begin{ejemplo}[Partícula que parte del reposo]

Un vehículo parte del reposo y mantiene una aceleración constante

\[ a=\SI{2.0}{\meter\per\second\squared}. \]

¿Cuánto se desplaza durante los primeros

\[ t=\SI{5.0}{\second}? \]

Como parte del reposo,

\[ v_0=0. \]

Utilizamos

\[ \Delta x=v_0t+\frac{1}{2}at^2. \]

Entonces,

\[
\Delta x
=
0
+
\frac{1}{2}
\left(\SI{2.0}{\meter\per\second\squared}\right)
\left(\SI{5.0}{\second}\right)^2.
\]

Por tanto,

\[
\Delta x
=
\frac{1}{2}(2.0)(25)\,\si{\meter}.
\]

Así,

\[ \boxed{\Delta x=\SI{25}{\meter}}. \]
\end{ejemplo}


\section{Gráfica posición-tiempo en MRUA}

La ecuación

\[ x=x_0+v_0t+\frac{1}{2}at^2 \]

es una función cuadrática del tiempo.

Por ello, la gráfica posición-tiempo es una parábola.

Si

\[ a>0, \]

la gráfica tiene concavidad hacia arriba.

Si

\[ a<0, \]

la gráfica tiene concavidad hacia abajo.

\begin{figure}[H]
    \centering

    \begin{tikzpicture}
        \begin{axis}[
            axis lines=middle,
            xlabel={$t$},
            ylabel={$x$},
            xmin=0,
            xmax=5,
            ymin=0,
            ymax=18,
            width=9cm,
            height=6cm
        ]

            \addplot[
                domain=0:4,
                samples=100,
                thick
            ] {1+x+0.5*x^2};

        \end{axis}
    \end{tikzpicture}

    \caption{Ejemplo de una gráfica posición-tiempo con aceleración constante
    positiva.}
\end{figure}

\begin{interpretacion}
La pendiente de la gráfica \(x-t\) continúa representando la velocidad.

Como la pendiente cambia con el tiempo, la velocidad también cambia.

La curvatura de la gráfica revela la presencia de aceleración.
\end{interpretacion}


\section{Tercera ecuación: relación entre velocidad y posición}

En algunos problemas conocemos la aceleración y el desplazamiento, pero no
conocemos el tiempo.

Conviene obtener una ecuación que no contenga \(t\).

Partimos de

\[ v=v_0+at. \]

Despejamos:

\[ t=\frac{v-v_0}{a}. \]

También tenemos

\[ \Delta x=\frac{v_0+v}{2}\,t. \]

Sustituimos el tiempo:

\[
\Delta x
=
\frac{v_0+v}{2}
\left(
\frac{v-v_0}{a}
\right).
\]

Utilizamos la identidad

\[
(v_0+v)(v-v_0)
=
v^2-v_0^2.
\]

Entonces,

\[
\Delta x
=
\frac{v^2-v_0^2}{2a}.
\]

Por tanto,

\[ \boxed{v^2=v_0^2+2a\Delta x}. \]

O, utilizando posiciones,

\[ \boxed{v^2=v_0^2+2a(x-x_0)}. \]

\begin{idea}
Esta ecuación es especialmente útil cuando el problema no proporciona el
tiempo ni lo solicita.
\end{idea}


\begin{ejemplo}[Frenado de un automóvil]

Un automóvil se mueve inicialmente con

\[ v_0=\SI{20}{\meter\per\second} \]

y frena con aceleración constante

\[ a=-\SI{4.0}{\meter\per\second\squared}. \]

Determine la distancia necesaria para detenerse.

Elegimos como positivo el sentido inicial del movimiento.

Al detenerse,

\[ v=0. \]

Utilizamos

\[ v^2=v_0^2+2a\Delta x. \]

Entonces,

\[
0
=
\left(\SI{20}{\meter\per\second}\right)^2
+
2
\left(-\SI{4.0}{\meter\per\second\squared}\right)
\Delta x.
\]

Por tanto,

\[
0
=
400\,\si{\meter\squared\per\second\squared}
-
8.0\,\si{\meter\per\second\squared}\,\Delta x.
\]

Despejando,

\[
\Delta x
=
\frac{400}{8.0}\,\si{\meter}.
\]

Así,

\[ \boxed{\Delta x=\SI{50}{\meter}}. \]

\begin{interpretacion}
La aceleración es negativa porque hemos elegido como positivo el sentido
inicial del movimiento.

Como velocidad y aceleración tienen signos opuestos, la rapidez disminuye
hasta llegar a cero.
\end{interpretacion}
\end{ejemplo}


\section{Las ecuaciones del MRUA}

Para aceleración constante y tomando \(t_0=0\), las relaciones principales
son:

\[
\boxed{v=v_0+at}
\]

\[
\boxed{\Delta x=\frac{v_0+v}{2}\,t}
\]

\[
\boxed{\Delta x=v_0t+\frac{1}{2}at^2}
\]

\[
\boxed{v^2=v_0^2+2a\Delta x}
\]

y

\[
\boxed{x=x_0+v_0t+\frac{1}{2}at^2}.
\]

\begin{advertencia}
Estas ecuaciones no constituyen cinco leyes físicas diferentes.

Son relaciones matemáticas equivalentes que se obtienen a partir de una
misma hipótesis física:

\[
\boxed{a=\text{constante}}.
\]

No deben utilizarse cuando la aceleración varía significativamente durante
el intervalo considerado.
\end{advertencia}


\section{¿Qué ecuación conviene utilizar?}

No es necesario memorizar una fórmula distinta para cada tipo de problema.

Una estrategia más útil consiste en identificar:

\begin{itemize}
    \item posición inicial \(x_0\);
    \item posición final \(x\);
    \item desplazamiento \(\Delta x\);
    \item velocidad inicial \(v_0\);
    \item velocidad final \(v\);
    \item aceleración \(a\);
    \item tiempo \(t\).
\end{itemize}

Después elegimos una ecuación que contenga los datos conocidos y la
incógnita, evitando introducir variables adicionales innecesarias.

\begin{table}[H]
    \centering

    \caption{Ecuaciones de MRUA y variable que no aparece en cada una.}

    \begin{tabular}{lc}
        \toprule
        Ecuación & Variable ausente \\
        \midrule
        \(v=v_0+at\) & \(\Delta x\) \\
        \(\Delta x=\dfrac{v_0+v}{2}t\) & \(a\) \\
        \(\Delta x=v_0t+\dfrac{1}{2}at^2\) & \(v\) \\
        \(v^2=v_0^2+2a\Delta x\) & \(t\) \\
        \bottomrule
    \end{tabular}
\end{table}

\begin{notapedagogica}
Este cuadro es más útil si se presenta después de haber deducido las
ecuaciones.

La intención no es que el estudiante escoja fórmulas mecánicamente, sino que
comprenda qué magnitudes relaciona cada una.
\end{notapedagogica}


\section{Gráfica aceleración-tiempo}

En MRUA,

\[ a=\text{constante}. \]

Por tanto, la gráfica \(a-t\) es una línea horizontal.

\begin{figure}[H]
    \centering

    \begin{tikzpicture}
        \begin{axis}[
            axis lines=middle,
            xlabel={$t$},
            ylabel={$a$},
            xmin=0,
            xmax=6,
            ymin=0,
            ymax=5,
            width=9cm,
            height=6cm
        ]

            \addplot[
                domain=0:5,
                samples=2,
                thick
            ] {2};

        \end{axis}
    \end{tikzpicture}

    \caption{Gráfica aceleración-tiempo para una aceleración constante
    positiva.}
\end{figure}

El área bajo la gráfica es

\[ A=a\Delta t. \]

Pero

\[ \Delta v=a\Delta t. \]

Por tanto,

\[ \boxed{\Delta v=A}. \]

\begin{interpretacion}
En una gráfica aceleración-tiempo, el área algebraica representa el cambio
de velocidad.
\end{interpretacion}


\section{Relación entre las tres gráficas}

Las gráficas de posición, velocidad y aceleración no son representaciones
independientes.

Están conectadas entre sí.

En una gráfica \(x-t\):

\[ \boxed{\text{pendiente}=v}. \]

En una gráfica \(v-t\):

\[ \boxed{\text{pendiente}=a}. \]

Además,

\[
\boxed{
\text{área bajo }v-t
=
\Delta x
}
\]

y

\[
\boxed{
\text{área bajo }a-t
=
\Delta v
}.
\]

Para MRUA:

\begin{itemize}
    \item \(a-t\) es una línea horizontal;
    \item \(v-t\) es una línea recta;
    \item \(x-t\) es una parábola.
\end{itemize}

\begin{idea}
Puede recordarse la cadena conceptual:

\[
x
\xrightarrow{\text{pendiente}}
v
\xrightarrow{\text{pendiente}}
a
\]

y, en sentido inverso,

\[
a
\xrightarrow{\text{área}}
\Delta v
\xrightarrow{\text{área}}
\Delta x.
\]
\end{idea}


\section{Cambio de sentido en un MRUA}

Una partícula puede reducir su rapidez, detenerse y posteriormente comenzar
a moverse en sentido contrario.

El cambio de sentido ocurre cuando

\[ v=0. \]

Si

\[ v=v_0+at, \]

podemos encontrar el instante de cambio de sentido haciendo

\[ 0=v_0+at. \]

Por tanto,

\[ \boxed{t=-\frac{v_0}{a}}. \]

Para que este instante sea físicamente posterior al inicial debe resultar

\[ t>0. \]

Esto ocurre cuando \(v_0\) y \(a\) tienen signos opuestos.


\begin{ejemplo}[Partícula que se detiene y regresa]

Una partícula tiene inicialmente

\[ v_0=+\SI{12}{\meter\per\second} \]

y aceleración constante

\[ a=-\SI{3}{\meter\per\second\squared}. \]

Determine cuándo se detiene momentáneamente.

Utilizamos

\[ v=v_0+at. \]

En el instante de detención,

\[ v=0. \]

Entonces,

\[
0
=
\SI{12}{\meter\per\second}
-
\left(\SI{3}{\meter\per\second\squared}\right)t.
\]

Por tanto,

\[ \boxed{t=\SI{4}{\second}}. \]

Antes de ese instante,

\[ v>0. \]

Después de ese instante,

\[ v<0. \]

Por ello, la partícula cambia de sentido.

\begin{interpretacion}
En \(t=\SI{4}{\second}\), la velocidad es cero, pero la aceleración continúa
siendo

\[ a=-\SI{3}{\meter\per\second\squared}. \]

Por eso la partícula no permanece detenida.
\end{interpretacion}
\end{ejemplo}


\section{Ejemplo integrador}

\begin{ejemplo}[Automóvil con aceleración constante]

Un automóvil se encuentra inicialmente en

\[ x_0=\SI{10}{\meter}. \]

Su velocidad inicial es

\[ v_0=\SI{4}{\meter\per\second} \]

y experimenta una aceleración constante

\[ a=\SI{2}{\meter\per\second\squared}. \]

Determine, después de

\[ t=\SI{5}{\second}: \]

\begin{enumerate}[label=\alph*)]
    \item la velocidad;
    \item el desplazamiento;
    \item la posición final.
\end{enumerate}
\end{ejemplo}

Para la velocidad utilizamos

\[ v=v_0+at. \]

Entonces,

\[
v
=
\SI{4}{\meter\per\second}
+
\left(\SI{2}{\meter\per\second\squared}\right)
\left(\SI{5}{\second}\right).
\]

Así,

\[ \boxed{v=\SI{14}{\meter\per\second}}. \]

Para el desplazamiento,

\[ \Delta x=v_0t+\frac{1}{2}at^2. \]

Sustituyendo,

\[
\Delta x
=
\left(\SI{4}{\meter\per\second}\right)
\left(\SI{5}{\second}\right)
+
\frac{1}{2}
\left(\SI{2}{\meter\per\second\squared}\right)
\left(\SI{5}{\second}\right)^2.
\]

Entonces,

\[
\Delta x
=
\SI{20}{\meter}
+
\SI{25}{\meter}.
\]

Por tanto,

\[ \boxed{\Delta x=\SI{45}{\meter}}. \]

Finalmente,

\[
x=x_0+\Delta x.
\]

Así,

\[
x
=
\SI{10}{\meter}
+
\SI{45}{\meter}.
\]

Por tanto,

\[ \boxed{x=\SI{55}{\meter}}. \]

\begin{interpretacion}
El automóvil comenzó a \(\SI{10}{\meter}\) del origen.

Durante los cinco segundos avanzó \(\SI{45}{\meter}\), mientras su velocidad
aumentó de \(\SI{4}{\meter\per\second}\) a
\(\SI{14}{\meter\per\second}\).
\end{interpretacion}


\section{Procedimiento para resolver problemas de MRUA}

\begin{notapedagogica}
Una estrategia útil para enseñar la resolución de problemas es:

\begin{enumerate}

    \item Dibujar un esquema.

    \item Elegir el eje y su sentido positivo.

    \item Identificar la posición inicial.

    \item Asignar signos a \(v_0\), \(v\), \(a\) y \(\Delta x\).

    \item Escribir los datos con unidades coherentes.

    \item Identificar la incógnita.

    \item Confirmar que el problema puede modelarse con aceleración constante.

    \item Elegir la relación que conecta los datos con la incógnita.

    \item Despejar algebraicamente antes de sustituir, cuando ayude a la
    claridad.

    \item Sustituir con unidades.

    \item Verificar dimensiones.

    \item Interpretar el signo y el valor del resultado.

\end{enumerate}

La parte más importante suele ocurrir antes de utilizar las ecuaciones:
\textbf{definir correctamente el sistema de referencia y los signos}.
\end{notapedagogica}


\section{Casos especiales importantes}

\subsection{Partícula que parte del reposo}

Si

\[ v_0=0, \]

entonces

\[ v=at, \]

\[ \Delta x=\frac{1}{2}at^2, \]

y

\[ v^2=2a\Delta x. \]

Estas ecuaciones son casos particulares de las ecuaciones generales, no
fórmulas nuevas.


\subsection{Aceleración nula}

Si

\[ a=0, \]

entonces

\[ v=v_0=\text{constante}. \]

Además,

\[ \Delta x=v_0t. \]

Recuperamos así el movimiento rectilíneo uniforme.

\begin{idea}
El MRU puede verse como el caso particular del movimiento con aceleración
constante en el que

\[ a=0. \]
\end{idea}


\subsection{Objeto que se detiene}

Cuando un móvil se detiene instantáneamente,

\[ v=0. \]

Esto no significa necesariamente que:

\[ a=0. \]

Si la aceleración continúa actuando, el móvil puede invertir posteriormente
su sentido.


\section{Condiciones de validez del modelo}

\begin{condicion}
Las ecuaciones de MRUA pueden utilizarse cuando:

\begin{itemize}
    \item el movimiento se analiza en una dimensión;
    \item la aceleración permanece constante durante el intervalo;
    \item se utiliza un único sistema de referencia;
    \item los signos de posiciones, velocidades y aceleraciones son
    consistentes con el eje elegido.
\end{itemize}
\end{condicion}

Si la aceleración cambia con el tiempo, las ecuaciones

\[ v=v_0+at \]

y

\[ x=x_0+v_0t+\frac{1}{2}at^2 \]

ya no describen, en general, el movimiento completo.


\section{Errores y confusiones frecuentes}

\begin{errorfrecuente}
\textbf{Creer que aceleración significa siempre aumento de rapidez.}

Un cuerpo también tiene aceleración cuando disminuye su rapidez o cambia de
sentido.
\end{errorfrecuente}


\begin{errorfrecuente}
\textbf{Pensar que una aceleración negativa significa frenado.}

El signo indica dirección.

La rapidez disminuye únicamente cuando velocidad y aceleración tienen signos
opuestos.
\end{errorfrecuente}


\begin{errorfrecuente}
\textbf{Confundir velocidad cero con aceleración cero.}

Un cuerpo puede detenerse instantáneamente y conservar una aceleración
distinta de cero.
\end{errorfrecuente}


\begin{errorfrecuente}
\textbf{Interpretar \(\si{\meter\per\second\squared}\) como una distancia.}

La aceleración expresa cambio de velocidad por unidad de tiempo.
\end{errorfrecuente}


\begin{errorfrecuente}
\textbf{Utilizar}

\[ \bar{v}=\frac{v_0+v}{2} \]

para cualquier movimiento.

Esta relación requiere aceleración constante.
\end{errorfrecuente}


\begin{errorfrecuente}
\textbf{Usar las ecuaciones de MRUA cuando la aceleración no es constante.}

Las ecuaciones cinemáticas deducidas en este capítulo dependen de la
hipótesis

\[ a=\text{constante}. \]
\end{errorfrecuente}


\begin{errorfrecuente}
\textbf{Escoger signos diferentes para cada ecuación.}

El sentido positivo debe establecerse una sola vez y mantenerse durante todo
el problema.
\end{errorfrecuente}


\begin{errorfrecuente}
\textbf{Asignar el signo de la aceleración basándose únicamente en que el
objeto ``frena''.}

Primero se define el eje.

Después se determina la dirección del vector aceleración.

El signo resulta de esa elección.
\end{errorfrecuente}


\begin{errorfrecuente}
\textbf{Confundir el área bajo una gráfica \(v-t\) con la distancia
recorrida.}

El área algebraica representa el desplazamiento.

Si la velocidad cambia de signo, para obtener la distancia deben sumarse las
magnitudes de los desplazamientos de cada tramo.
\end{errorfrecuente}


\section{Secuencia conceptual para explicar el tema}

\begin{notapedagogica}
Una secuencia clara para dictar esta clase es:

\begin{enumerate}

    \item Recordar que la velocidad describe el cambio de posición.

    \item Plantear situaciones en las que la velocidad cambia.

    \item Preguntar qué magnitud necesitamos para describir ese cambio.

    \item Definir aceleración media:

    \[
    \bar{a}=\frac{\Delta v}{\Delta t}.
    \]

    \item Explicar cuidadosamente la unidad

    \[
    \si{\meter\per\second\squared}.
    \]

    \item Resolver un ejemplo sencillo de aceleración media.

    \item Introducir aceleración instantánea mediante intervalos cada vez más
    pequeños.

    \item Presentar

    \[
    a=\frac{\dd v}{\dd t}.
    \]

    \item Comparar los signos de \(v\) y \(a\).

    \item Establecer que misma señal significa aumento de rapidez y signos
    opuestos disminución de rapidez.

    \item Mostrar ejemplos con \(v<0\) para evitar que el estudiante asocie
    automáticamente aceleración negativa con frenado.

    \item Analizar la pendiente de una gráfica \(v-t\).

    \item Introducir el MRUA como el caso de aceleración constante.

    \item Construir primero

    \[
    v=v_0+at.
    \]

    \item Utilizar la gráfica \(v-t\) para obtener el desplazamiento.

    \item Introducir

    \[
    \bar{v}=\frac{v_0+v}{2}
    \]

    únicamente bajo la condición de aceleración constante.

    \item Deducir

    \[
    \Delta x=v_0t+\frac{1}{2}at^2.
    \]

    \item Relacionar esta ecuación con la gráfica parabólica \(x-t\).

    \item Eliminar el tiempo para obtener

    \[
    v^2=v_0^2+2a\Delta x.
    \]

    \item Relacionar las tres gráficas \(x-t\), \(v-t\) y \(a-t\).

    \item Resolver un ejemplo de frenado.

    \item Resolver un ejemplo donde el móvil se detiene y cambia de sentido.

    \item Finalizar mostrando que todas las ecuaciones provienen de una sola
    condición:

    \[
    a=\text{constante}.
    \]

\end{enumerate}
\end{notapedagogica}


\section{Preguntas para comprobar comprensión}

Antes de pasar a una colección formal de ejercicios pueden plantearse
preguntas como:

\begin{enumerate}

    \item ¿Qué significa físicamente que un cuerpo tenga aceleración?

    \item ¿Qué significa una aceleración de

    \[
    \SI{4}{\meter\per\second\squared}?
    \]

    \item ¿Puede un cuerpo tener aceleración aunque su rapidez disminuya?

    \item ¿Una aceleración negativa significa necesariamente que el cuerpo
    está frenando?

    \item Si \(v>0\) y \(a<0\), ¿la rapidez aumenta o disminuye?

    \item Si \(v<0\) y \(a<0\), ¿la rapidez aumenta o disminuye?

    \item ¿Puede ocurrir simultáneamente

    \[
    v=0
    \]

    y

    \[
    a\neq0?
    \]

    \item ¿Qué representa la pendiente de una gráfica \(v-t\)?

    \item ¿Qué representa el área bajo una gráfica \(v-t\)?

    \item ¿Qué representa el área bajo una gráfica \(a-t\)?

    \item ¿Qué forma tiene la gráfica \(v-t\) cuando la aceleración es
    constante?

    \item ¿Qué forma tiene la gráfica \(x-t\) en MRUA?

    \item ¿Por qué la relación

    \[
    \bar{v}=\frac{v_0+v}{2}
    \]

    requiere aceleración constante?

    \item ¿Qué ecuación de MRUA resulta especialmente conveniente si no
    conocemos el tiempo?

    \item Si \(a=0\), ¿qué tipo de movimiento obtenemos?

    \item ¿Qué condición permite identificar un instante de cambio de
    sentido?

\end{enumerate}


\section{Conexión con el tema siguiente}

Las ecuaciones desarrolladas en este capítulo permiten describir cualquier
movimiento rectilíneo cuya aceleración sea constante.

Existe un caso particularmente importante en la naturaleza: el movimiento
vertical de un objeto cuando la única interacción relevante es la gravedad
y podemos despreciar la resistencia del aire.

Cerca de la superficie terrestre, este movimiento puede modelarse mediante
una aceleración aproximadamente constante.

En el siguiente capítulo aplicaremos por tanto toda la estructura del MRUA
al estudio de la \textbf{caída libre y el movimiento vertical}.


\begin{resumen}

\begin{itemize}

    \item La aceleración describe el cambio de velocidad con respecto al
    tiempo.

    \item La aceleración media es

    \[
    \boxed{\bar{a}=\frac{\Delta v}{\Delta t}}.
    \]

    \item La aceleración instantánea es

    \[
    \boxed{a=\frac{\dd v}{\dd t}}
    \]

    y también

    \[
    \boxed{a=\frac{\dd^2x}{\dd t^2}}.
    \]

    \item La unidad SI de aceleración es

    \[
    \si{\meter\per\second\squared}.
    \]

    \item Si \(v\) y \(a\) tienen el mismo signo, la rapidez aumenta.

    \item Si \(v\) y \(a\) tienen signos opuestos, la rapidez disminuye.

    \item Una aceleración negativa no implica necesariamente disminución de
    rapidez.

    \item Una velocidad instantánea nula no implica necesariamente
    aceleración nula.

    \item En una gráfica \(v-t\), la pendiente representa la aceleración.

    \item En MRUA,

    \[
    \boxed{a=\text{constante}}.
    \]

    \item La velocidad cambia linealmente con el tiempo:

    \[
    \boxed{v=v_0+at}.
    \]

    \item Para aceleración constante,

    \[
    \boxed{\bar{v}=\frac{v_0+v}{2}}.
    \]

    \item El desplazamiento puede calcularse mediante

    \[
    \boxed{\Delta x=\frac{v_0+v}{2}\,t}.
    \]

    \item También,

    \[
    \boxed{\Delta x=v_0t+\frac{1}{2}at^2}.
    \]

    \item La posición está dada por

    \[
    \boxed{x=x_0+v_0t+\frac{1}{2}at^2}.
    \]

    \item La relación que no contiene explícitamente el tiempo es

    \[
    \boxed{v^2=v_0^2+2a\Delta x}.
    \]

    \item En MRUA, la gráfica \(a-t\) es una línea horizontal.

    \item La gráfica \(v-t\) es una línea recta.

    \item La gráfica \(x-t\) es una parábola.

    \item El área bajo una gráfica \(v-t\) representa desplazamiento.

    \item El área bajo una gráfica \(a-t\) representa cambio de velocidad.

    \item Todas las ecuaciones cinemáticas del MRUA dependen de la condición

    \[
    a=\text{constante}.
    \]

\end{itemize}

\end{resumen}